\chapter{Contexte, problématique et objectifs}
\label{ch:contexte}

\section{Contexte du projet}

Dans la plupart des organisations, les données de gestion sont stockées dans des bases relationnelles
transactionnelles (OLTP). Ces bases sont optimisées pour enregistrer rapidement de nombreuses petites
opérations --- une commande, une facture, un paiement --- et sont donc fortement normalisées : une simple
question comme « quel est le chiffre d'affaires par catégorie de client et par année ? » nécessite des
jointures entre une demi-douzaine de tables, des règles de calcul (hors taxes, toutes taxes comprises,
avoirs) et une bonne connaissance du schéma.

L'informatique décisionnelle répond traditionnellement à ce problème par la construction d'un entrepôt de
données (\emph{data warehouse}) alimenté par des traitements ETL et exploité au moyen de rapports et de
tableaux de bord. Cette approche reste cependant rigide : chaque nouveau besoin d'analyse passe par
l'équipe technique, qui doit écrire la requête ou faire évoluer le rapport.

\section{Problématique}

Le sujet formule la problématique suivante : \emph{comment concevoir une plateforme de données moderne
capable d'automatiser l'ingestion, la transformation et la modélisation des données, tout en offrant une
interface conversationnelle permettant d'interroger ces données en langage naturel grâce à
l'intelligence artificielle ?}

Cette problématique soulève plusieurs questions techniques :
\begin{itemize}
  \item comment extraire de manière fiable et \textbf{incrémentale} les données d'un système source
        transactionnel vers un entrepôt dans le cloud ?
  \item comment organiser les données en couches successives garantissant la \textbf{traçabilité}, la
        \textbf{qualité} et la \textbf{facilité d'analyse} ?
  \item comment automatiser, surveiller et fiabiliser l'enchaînement des traitements ?
  \item comment permettre à un LLM de générer un SQL \textbf{correct} sur un schéma métier, et comment
        garantir que cette génération reste \textbf{sûre} (aucune modification possible des données) ?
\end{itemize}

\section{Objectifs}

\subsection{Objectif général}
Concevoir et mettre en œuvre une Modern Data Platform centralisant les données de l'entreprise, les
transformant selon une architecture Médaillon, et exposant une application Text-to-SQL capable de
convertir automatiquement les questions des utilisateurs en requêtes SQL.

\subsection{Objectifs spécifiques}
\begin{enumerate}
  \item Concevoir une architecture moderne de données de bout en bout.
  \item Automatiser les pipelines d'ingestion de SQL Server vers Snowflake avec DLT.
  \item Mettre en place les couches Bronze, Silver et Gold dans Snowflake.
  \item Transformer, documenter et tester les données avec dbt.
  \item Orchestrer l'ensemble des traitements avec Dagster.
  \item Développer une interface utilisateur avec Streamlit.
  \item Intégrer Snowflake Cortex AI pour convertir les questions en langage naturel en requêtes SQL
        exécutables.
\end{enumerate}

Au-delà de ces objectifs, le projet traite les pistes d'amélioration identifiées en cours de réalisation :
historisation des données, modèles incrémentaux, déploiement dans Streamlit in Snowflake, alertes et
intégration continue.

\section{Le jeu de données Wide World Importers}

\emph{Wide World Importers} (WWI) est la base d'exemple de référence publiée par Microsoft pour SQL
Server depuis la version 2016. Elle modélise l'activité d'une entreprise fictive de distribution en gros
d'articles de nouveauté (jouets, mugs, t-shirts, articles d'emballage, gadgets USB\ldots) basée aux
États-Unis, entre janvier 2013 et mai 2016.

Microsoft fournit deux sauvegardes :
\begin{itemize}
  \item \texttt{WideWorldImporters-Full.bak} : la base \textbf{transactionnelle (OLTP)} ;
  \item \texttt{WideWorldImportersDW-Full.bak} : un entrepôt de données déjà modélisé (OLAP).
\end{itemize}

Le projet utilise volontairement la base \textbf{OLTP} : elle joue le rôle du système source, et c'est
la plateforme développée qui reconstruit l'entrepôt analytique. Utiliser la version DW aurait vidé le
travail de modélisation de sa substance.

La base comprend quatre schémas métier, présentés dans le tableau~\ref{tab:wwi-schemas}.

\begin{table}[H]
\centering
\caption{Schémas métier de la base WideWorldImporters (OLTP)}
\label{tab:wwi-schemas}
\begin{tabularx}{\linewidth}{@{}l X r@{}}
\toprule
\textbf{Schéma} & \textbf{Contenu et principales tables} & \textbf{Lignes (max.)} \\
\midrule
\texttt{Sales} & Clients, commandes, factures, transactions clients (\texttt{Invoices},
  \texttt{InvoiceLines}, \texttt{Orders}, \texttt{OrderLines}, \texttt{Customers}) & 231\,412 \\
\texttt{Purchasing} & Fournisseurs, bons de commande, transactions fournisseurs & 8\,367 \\
\texttt{Warehouse} & Articles, stocks, mouvements de stock, groupes d'articles & 236\,667 \\
\texttt{Application} & Référentiels : villes, provinces, pays, personnes, modes de livraison et de
  paiement & 37\,940 \\
\bottomrule
\end{tabularx}
\end{table}

Plusieurs caractéristiques de WWI en font un cas d'étude réaliste et exigeant :
\begin{itemize}
  \item des \textbf{tables temporelles} (\emph{system-versioned}) qui conservent l'historique de chaque
        modification dans des tables d'archive (\texttt{*\_Archive}) ;
  \item une colonne \texttt{LastEditedWhen} sur les tables transactionnelles, exploitable pour
        l'ingestion incrémentale ;
  \item des types de données spécifiques à SQL Server (\texttt{geography}, \texttt{varbinary}) et des
        champs JSON (\texttt{CustomFields}) ;
  \item des tables \emph{In-Memory OLTP} (optimisées en mémoire).
\end{itemize}
